\VerdierPhysicalPageAnchor{0071}%
\VerdierSemanticAnchor{ch1-2-3}
\subsection{Graded structure on the categories of complexes}

\VerdierSemanticAnchor{ch1-2-3-1}
\noindent\textbf{2.3.1.}\quad
Let $Y$ be an $n$-fold complex and let $i\in[n]$. Set:

\VerdierSemanticAnchor{ch1-2-3-1-eq1}
\begin{equation}
  \begin{aligned}
    (T_iY)^\sigma&=Y^{\sigma+e_i},
      &&\sigma\in\mathbb Z[n],\\
    d_{T_iY}^{\sigma,j}&=\nu(i,j)d_Y^{\sigma+e_i,j},
      &&j\in[n],
  \end{aligned}
  \tag{2.3.1.1}
\end{equation}
where $\nu(i,j)=1$ if $i\ne j$, and $\nu(i,j)=-1$ if $i=j$.

\VerdierSemanticAnchor{ch1-2-3-1-morphism-shift}
For every morphism of $n$-fold complexes
$f\colon Y\longrightarrow Y'$, set:

\VerdierSemanticAnchor{ch1-2-3-1-eq2}
\begin{equation}
  (T_if)^\sigma=f^{\sigma+e_i},
  \qquad \sigma\in\mathbb Z[n].
  \tag{2.3.1.2}
\end{equation}

\VerdierSemanticAnchor{ch1-2-3-1-translation-functor}
This defines a functor
$T_i\colon\operatorname{comp}^{n}(C)\longrightarrow
\operatorname{comp}^{n}(C)$, called the
\emph{translation functor in direction $i$}.

\VerdierSemanticAnchor{ch1-2-3-1-n-equals-one}
When $n=1$, we use the notation $TY$, $Tf$, or else $Y[1]$, $f[1]$.
The translation functor in direction $i$ is an automorphism of
$\operatorname{comp}^{n}(C)$.

\VerdierSemanticAnchor{ch1-2-3-1-eq3}
The functors $T_i$, $i\in[n]$, commute pairwise:
\begin{equation}
  T_iT_j=T_jT_i,
  \qquad \forall i,j\in[n].
  \tag{2.3.1.3}
\end{equation}

\VerdierSemanticAnchor{ch1-2-3-1-action}
We can therefore define an action of the group $\mathbb Z[n]$ on the category
$\operatorname{comp}^{n}(C)$:

\VerdierSemanticAnchor{ch1-2-3-1-eq4}
\begin{equation}
  \sigma=\sum_{i\in[n]}l_ie_i,
  \qquad
  T(\sigma)=T_1^{l_1}\cdots T_n^{l_n}.
  \tag{2.3.1.4}
\end{equation}

\VerdierSemanticAnchor{ch1-2-3-1-strict-graded-category}
Henceforth the notation $\operatorname{comp}^{n}(C)$ will denote the category
$\operatorname{comp}^{n}(C)$ endowed with the strict $\mathbb Z[n]$-category
structure described above.

\VerdierSemanticAnchor{ch1-2-3-2}
\noindent\textbf{2.3.2.}\quad
Let $\varphi\colon[n]\longrightarrow[m]$ be a map, let $Y$ be an $n$-fold
complex in $C$, and let $i$ be an element of $[n]$. We suppose, to fix ideas,
that countable sums are representable in $C$.

\VerdierSemanticAnchor{ch1-2-3-2-equal-objects}
First observe that the objects
$\left(\int_\varphi^\Sigma T_iY\right)^\tau$ and
$\left(T_{\varphi(i)}\int_\varphi^\Sigma Y\right)^\tau$ are equal.
Indeed, by definition they are both equal to:

\VerdierSemanticAnchor{ch1-2-3-2-direct-sum}
\begin{equation*}
  \bigoplus_{\varphi(\sigma)=\tau+e_{\varphi(i)}}Y^\sigma.
\end{equation*}
Now set:

\VerdierSemanticAnchor{ch1-2-3-2-eq1}
\begin{equation}
  \begin{aligned}
    \bigl(p_{i,\varphi}(Y)\bigr)^\tau\colon
    \left(\int_\varphi^\Sigma T_iY\right)^\tau
      &\longrightarrow
    \left(T_{\varphi(i)}\int_\varphi^\Sigma Y\right)^\tau,
      &&\tau\in\mathbb Z[m],\\[2pt]
    \bigl(p_{i,\varphi}(Y)\bigr)^\tau
      &=\bigoplus_{\varphi(\sigma)=\tau+e_{\varphi(i)}}
        \varepsilon(\sigma,i)\varepsilon(\varphi(\sigma),\varphi(i))
        \operatorname{id}_{Y^\sigma}.
  \end{aligned}
  \tag{2.3.2.1}
\end{equation}
